\documentclass[aspectratio=169]{beamer}
\usepackage{etex} % fixes new-dimension error
\usepackage{lmodern}
\usepackage[T1]{fontenc}
\usepackage{graphicx,amsmath}
\usepackage{stmaryrd} % cf. interleave
\usepackage{booktabs}
\usepackage{amscd}
\usepackage{multicol}
\usepackage[absolute,overlay]{textpos}
\usepackage{alltt}
\usepackage{proof}
\usepackage{subcaption}
\usepackage{tikz}
\usepackage{tikz-cd}
\usepackage[new]{old-arrows}
\usepackage[all]{xy}
\usepackage{pgfplots}
\usepackage{textcomp}
\usepackage{cancel}

% Cyrillic characters -- \foreignlanguage{russian}{слово} OR \begin{otherlanguage*}{russian}... 
\usepackage[T2A,T1]{fontenc}
\usepackage[utf8]{inputenc}
\usepackage[russian,english]{babel}


\usepackage{transparent}
\usepackage{xspace}
\usepackage{listings}
\usepackage{pdfpages}
\usepackage{relsize}



%%%%%%%%%%%%% Macros
% \newcommand{\Ban}{\catfont{Ban}}
% \newcommand{\Cats}{\catfont{Cat}}

%%%% Misc
%% Operations
\newcommand{\sem}[1]{\llbracket #1 \rrbracket}
\DeclareMathOperator{\img}{\mathrm{im}}
\DeclareMathOperator{\dom}{\mathrm{dom}}
\DeclareMathOperator{\codom}{\mathrm{codom}}
%% Sets of numbers
\newcommand{\N}{\mathbb{N}}
\newcommand{\Z}{\mathbb{Z}}
\newcommand{\Nats}{\mathbb{N}}
\newcommand{\Reals}{\mathbb{R}}
\newcommand{\Rz}{\Reals_{\geq 0}}
\newcommand{\Complex}{\mathbb{C}}
%% Writing
\newcommand{\cf}{\emph{cf.}}
\newcommand{\ie}{\emph{i.e.}}
\newcommand{\eg}{\emph{e.g.}}
\newcommand{\df}[1]{\emph{\textbf{#1}}}
%%%%%%%%%%%%%%%% End of Misc

%%%% Programming Stuff
%% Types
\newcommand{\typefont}[1]{\mathbb{#1}}
\newcommand{\typeOne}{1}
\newcommand{\typeTwo}{2}
\newcommand{\typeA}{\typefont{A}}
\newcommand{\typeX}{\typefont{X}}
\newcommand{\typeB}{\typefont{B}}
\newcommand{\typeC}{\typefont{C}}
\newcommand{\typeV}{\typefont{V}}
\newcommand{\typeD}{\typefont{D}}
\newcommand{\typeI}{\typefont{I}}
%% RuleName
\newcommand{\rulename}[1]{(\mathrm{#1})}
%% Sequents
\newcommand{\jud}{\vdash}
\newcommand{\vljud}{\rhd}
\newcommand{\cojud}{\vdash_{\co}}
\newcommand{\vl}{\mathtt{v}}
\newcommand{\co}{\mathtt{c}}
% Program font
\newcommand{\prog}[1]{\mathtt{#1}}
\newcommand{\pseq}[3]{#1 \leftarrow #2; #3}
\newcommand{\ppm}[4]{(#1,#2) \leftarrow #3; #4}
\newcommand{\pinl}[1]{\prog{inl}(#1)}
\newcommand{\pinr}[1]{\prog{inr}(#1)}
\newcommand{\pcase}[4]{\prog{ case } #1 \prog{ of } \pinl{#2} \Rightarrow #3 ; \pinr{#2} \Rightarrow #4}
%% Sets of terms
\newcommand{\ValuesBP}[2]{\mathsf{Values}(#1, #2)}
\newcommand{\TermsBP}[2]{\mathsf{Terms}(#1, #2)}
\newcommand{\closValP}[1]{\ValuesBP{\emptyset}{#1}}
\newcommand{\closTermP}[1]{\TermsBP{\emptyset}{#1}}
\newcommand{\closVal}{\closValP{\typeA}}
\newcommand{\closTerm}{\closTermP{\typeA}}
%% Contextual equivalence
\newcommand{\ctxeq}{\equiv_{\prog{ctx}}}
%%%% End of Programming Stuff


% Misc by José
\newcommand{\wrap}[2][]{\begin{tabular}[#1]{@{}c@{}}#2\end{tabular}}
\newcommand{\mwrap}[1]{\ensuremath{\begin{array}{@{}c@{}}#1\end{array}}}
\def\trans#1{\xrightarrow{#1}}  % - a - > 
\def\Trans#1{\stackrel{#1}{\Longrightarrow}} % =a=> 
\newcommand{\transp}[2][35]{\color{fg!#1}#2}
\newcommand{\transpt}[2][.35]{\tikz{\node[inner sep=1pt,fill opacity=0.5]{#2}}}
\newcommand{\faded}[2][0.4]{{\transparent{#1}#2}} % alternative to "transp" using transparent package
\newcommand{\set}[1]{\left\{ #1 \right\}} % {a,b,...z}
\newcommand{\mi}[1]{\ensuremath{\mathit{#1}}\xspace}
\newcommand{\mf}[1]{\ensuremath{\mathsf{#1}}\xspace}
% \newcommand{\gold}[1]{\textcolor{darkgoldenrod}{#1}\xspace}


%------ using color ---------------------------------------------------------
\definecolor{goldenrod}{rgb}{.80392 .60784 .11373}
\definecolor{darkgoldenrod}{rgb}{.5451 .39608 .03137}
\definecolor{brown}{rgb}{.15 .15 .15}
\definecolor{darkolivegreen}{rgb}{.33333 .41961 .18431}
\definecolor{myGray}{gray}{0.85}
%
%
\newcommand{\red}[1]{\textcolor{red!80!black}{#1}\xspace}
\newcommand{\blue}[1]{\textcolor{blue}{#1}\xspace}
\newcommand{\gold}[1]{\textcolor{darkgoldenrod}{#1}\xspace}
\newcommand{\gray}[1]{\textcolor{myGray}{#1}\xspace}
% \def\alert#1{{\darkgoldenrod #1}}
% \def\alert#1{{\alert{#1}}}
%\def\brw#1{{\brown #1}}
% \def\structure#1{{\blue #1}}
% \def\tstructure#1{\textbf{\darkblue #1}}
%%\def\gre#1{{\green #1}}
\def\gre#1{{\darkolivegreen #1}}
\def\gry#1{{\textcolor{gray}{#1}}}
\def\rdb#1{{\red #1}}
\def\st{\mathbf{.}\,}
\def\laplace#1#2{*\txt{\mbox{ \fcolorbox{black}{myGray}{$\begin{array}{c}\mbox{#1}\\\\#2\\\\\end{array}$} }}}
%\newcommand{\galois}[2]{#1\; \dashv\; #2}



\def\ainv#1{\overline{#1}}
\def\aconv#1{#1^{\circ}} 
\def\rtran#1{\stackrel{#1}{\longrightarrow}}
\def\pair#1{\langle #1 \rangle}
\def\setdef#1#2{\mathopen{\{} #1 \asor #2 \mathclose{\}}}
\def\imp{\mathbin{\Rightarrow}}
\def\dimp{\mathbin{\Leftrightarrow}}
\def\rimp{\mathbin{\Leftarrow}}
\def\rra{\longrightarrow}
% \def\rcb#1#2#3#4{\def\nothing{}\def\range{#3}\mathopen{\langle}#1 \ #2 \ \ifx\range\nothing::\else: \ #3 :\fi \ #4\mathclose{\rangle}}
\def\abv{\stackrel{\rm abv}{=}}



% Spliting frames in 2 columns
\newcommand{\splittwo}[4]{ 
  \begin{columns}[T]% align columns
  \begin{column}{#1\textwidth} #3 \end{column} ~~~
  \begin{column}{#2\textwidth} #4 \end{column} \end{columns}
}\newcommand{\frsplit}[3][.48]{
  \begin{columns}%[T] % align columns
  \begin{column}{#1\textwidth} #2 \end{column} ~~~
  \begin{column}{#1\textwidth} #3 \end{column} \end{columns}
}
\newcommand{\frsplitdiff}[5][]{
  \begin{columns}[#1]%[T] % align columns
  \begin{column}{#2\textwidth} #4 \end{column} ~~~
  \begin{column}{#3\textwidth} #5 \end{column} \end{columns}
}
\newcommand{\frsplitt}[3][.48]{
  \begin{columns}[T] % align columns
  \begin{column}{#1\textwidth} #2 \end{column} ~~~
  \begin{column}{#1\textwidth} #3 \end{column} \end{columns}
}
\newcommand{\col}[2][.48]{\begin{column}{#1\textwidth} #2 \end{column}}
\newcommand{\colb}[3][.48]{\begin{column}{#1\textwidth} \begin{block}{#2} #3 \end{block} \end{column}}

% Spliting frames in 3 columns
\newcommand{\splitthree}[6]{
  \begin{columns}[T] % align columns
  \begin{column}{#1\textwidth} #4 \end{column} ~~~
  \begin{column}{#2\textwidth} #5 \end{column} ~~~
  \begin{column}{#3\textwidth} #6 \end{column} \end{columns}
}
\newcommand{\frsplitthree}[4][.31]{
  \begin{columns}%[T] % align columns
  \begin{column}{#1\textwidth} #2 \end{column} ~~~
  \begin{column}{#1\textwidth} #3 \end{column} ~~~
  \begin{column}{#1\textwidth} #4 \end{column} \end{columns}
}
\newcommand{\frsplitdiffthree}[7][]{
  \begin{columns}[#1]%[T] % align columns
  \begin{column}{#2\textwidth} #5 \end{column} ~~~
  \begin{column}{#3\textwidth} #6 \end{column} ~~~
  \begin{column}{#4\textwidth} #7 \end{column} \end{columns}
}
\newcommand{\frsplittthree}[4][.32]{
  \begin{columns}[T] % align columns
  \begin{column}{#1\textwidth} #2 \end{column} ~
  \begin{column}{#1\textwidth} #3 \end{column} ~
  \begin{column}{#1\textwidth} #4 \end{column} \end{columns}
}


\newcommand{\typerule}[4][]{\ensuremath{\begin{array}[#1]{c}\textsf{\scriptsize ({#2})} \\#3 \\\hline\raisebox{-3pt}{\ensuremath{#4}}\end{array}}}
\newcommand{\styperule}[3][]{\ensuremath{\begin{array}[#1]{c} #2 \\[0.5mm]\hline\raisebox{-4pt}{\ensuremath{#3}}\end{array}}}
\newcommand{\shrk}{\vspace{-3mm}}

\def\caixa#1{\medskip
  \begin{center}
  \fbox{\begin{minipage}{0.9\textwidth}\protect{#1}\end{minipage}}
  \end{center}}

\newcommand{\mybox}[2][4mm]{
  % \begin{minipage}{#1\textwidth}\begin{block}{}\centering #2\end{block}\end{minipage}}
  \tikz{\node[fill=barcolor!60,align=center,inner sep=#1]{#2};}}
\newcommand{\mycbox}[2][4mm]{
  \begin{center}\mybox[#1]{#2}\end{center}}
  % {\\[-5mm]\centering\mybox[#1]{#2}\\[-5mm]}}

%%%%% Tikz
% \usetikzlibrary{arrows.meta, calc, fit, tikzmark}
\usetikzlibrary{%
  positioning
 ,patterns
 ,arrows
 ,arrows.meta
 ,automata
 ,calc
 ,shapes
 ,fit
 ,tikzmark
 ,fadings
 ,decorations.pathreplacing
 ,plotmarks
% ,pgfplots.groupplots
 ,decorations.markings
 ,shadows
}
% \tikzset{shorten >=1pt,node distance=2cm,on grid,auto,initial text={},inner sep=2pt}
\tikzstyle{aut}=[shorten >=1pt,node distance=2cm,on grid,auto,initial text={},inner sep=2pt]
\tikzstyle{st}=[circle,draw=black,fill=black!10,inner sep=3pt]
\tikzstyle{sst}=[rectangle,draw=none,fill=none,inner sep=3pt]
\tikzstyle{final}=[accepting]

%%% Uppaal-like diagrams
\newcommand{\uppbox}[3][20mm]{\tikz{
  \node[black!15,fill=black!15,minimum width=#1,align=left](title){\textbf{{\footnotesize #2}}};
  \node[black!15,fill=black!15,left,xshift=4mm]at(title.east){\textbf{{\footnotesize #2}}};
  \node[blue!60!cyan,right] at(title.west){\textbf{{\footnotesize #2}}};
  \node[below,inner sep=2mm,fill=white,xshift=2mm](box)at(title.south){\includegraphics[width=#1]{#3}};
  \node[fit=(title)(box),draw=black,inner sep=0pt]{};
}}

\newcommand{\uppboxv}[3][20mm]{\tikz{
  \node[below,inner sep=2mm,fill=white](box){\includegraphics[height=#1]{#3}};
  \coordinate[yshift=5mm](top)at(box.north);
  \node[fit=(top)(box.north west)(box.north east),inner sep=0pt,fill=black!15](title){};
  \node[blue!60!cyan,right] at(title.west){\textbf{{\footnotesize #2}}};
  \node[fit=(title)(box),draw=black,inner sep=0pt]{};
}}

%%% frame for pictures [graphics-options]{content}
\newcommand{\includegraphicsframed}[2][]{\tikz{
  \node[below,inner sep=1mm,fill=barcolor,%xshift=2mm,
    draw=black,
    drop shadow={
      top color=gray,
      bottom color=white,
      %fill=gray,
      opacity=0.2,
      shadow xshift=4pt,
      shadow yshift=-3pt
    }](box){\includegraphics[#1]{#2}};
}}


%% COnfiguring Listings
\lstset{ % basic style
  language=scala,
  basicstyle=\ttfamily\scriptsize,
  breakatwhitespace=true,
  breaklines=true,
  % mathescape,
  % morecomment=[l]{//},
  % morecomment=[n]{/*}{*/},
  % frame=single,                    % adds a frame around the code
  rulecolor=\color{black!40},         % if not set, the frame-color may be changed on line-breaks within not-black text (e.g. comments (green here))
  xleftmargin=1.5mm,
  xrightmargin=1.5mm,
  backgroundcolor=\color{black!5},
  % line numbers
%  numbers=left,  % where to put the line-numbers; possible values are (none, left, right)
 numbersep=5pt, % how far the line-numbers are from the code
 numberstyle=\tiny\color{gray},   
 stepnumber=1,  % the step between two line-numbers. If it is 1 each line will be numbered      
%  xleftmargin=3mm,
%  xrightmargin=1.5mm,
%%%%%
  captionpos=b, % t or b (top or bottom)
  belowcaptionskip=5mm,
%%%%%
  % alsoletter={-},
  % emphstyle=\ttfamily\color{blue}, %\underbar,
  % emphstyle={[2]\ttfamily\color{green!50!black}},
  emphstyle=\bfseries\itshape\color{blue!80!black},       % moreemph={...} - layer keywords
  emphstyle={[2]\itshape\color{red!70!black}},%\underbar} % moreemph={[2]...} - inner keywords
  %
  keywordstyle=\bf\ttfamily\color{red!50!black},
  commentstyle=\sl\ttfamily\color{gray!70},
  % commentstyle=\color{green!60!black},
  stringstyle=\ttfamily\color{green!60!black},
  morestring=[b]",
  % morecomment=[l]{\#},
  frame=single,
  % numberstyle=\tiny, numbers=left, stepnumber=1, firstnumber=1, numberfirstline=true,
  %emph={act,proc,init,sort,map,var,eqn},
  %emph={[2]block,hide,comm,rename,allow,||,<>,sum,&&,=>,true,false},
  literate={\\§}{{{\mbox{\textdollar}}}}1,
  % literate=*{->}{{{\color{red!70!black}$\to$}}}{1}
  %            {.}{{{\color{red!70!black}.}}}{1}
  %            {+}{{{\color{red!70!black}\hspace*{1pt}+\hspace*{1pt}}}}{1}
  %            {|}{{{\color{red!70!black}|}}}{1}
  %            {||}{{{\color{red!70!black}|\!\!|}}}{1}
  %            {*}{{{\color{red!70!black}*}}}{1}
  %            {\#}{{{\color{red!70!black}\#}}}{1}
  %            {&}{{{\color{red!70!black}\&}}}{1}
  %            {:=}{{{\color{red!70!black}:=}}}{1}
  %            {=>}{{{\color{red!70!black}=>}}}{1}
  %            % {>}{{{\color{green!65!black}\hspace*{1pt}>\hspace*{1.5pt}}}}{2}
  %            % {<}{{{\color{green!65!black}\hspace*{1.5pt}<\hspace*{1pt}}}}{2}
  %            % {]}{{{\color{green!65!black}\hspace*{1pt}]\hspace*{1.5pt}}}}{1}
  %            % {[}{{{\color{green!65!black}\hspace*{1.5pt}[\hspace*{1pt}}}}{1}
  %            {>}{{{\color{green!65!black}>}}}{1}
  %            {<}{{{\color{green!65!black}<}}}{1}
  %            {]}{{{\color{green!65!black}]}}}{1}
  %            {[}{{{\color{green!65!black}[}}}{1}
  %  morekeywords={Merger1,Fifo2,Lossy3,Init1,Init2}
  % ,emph={act,proc,init,sort,map,var,eqn}
  % ,emph={[2]block,hide,comm,rename,allow,||,<>,sum,&&,=>}
}

\newcommand{\code}[1]{{\relsize{-1}\ttfamily #1}}
\newcommand{\cod}[1]{{\relsize{-0.5}\color{blue!75!black}\ttfamily #1}}
% \definecolor{mydarkgreen}{rgb}{0,0.6,0}
% \newcommand{\bash}[1]{\lstinline[basicstyle=\ttfamily\relsize{-0.5}\color{mydarkgreen},keywordstyle=\bf\sffamily\color{purple},columns=fullflexible,keepspaces,literate=*]|#1|}

\lstdefinestyle{tiny}{basicstyle=\ttfamily\relsize{-7}}

% Include slides from others
\newcommand{\byothers}[3]{{
\begin{frame}{}~\mycbox{\Large slides by #1\\pages #2}\end{frame}{}
\setbeamercolor{background canvas}{bg=}
\includepdf[pages=#2]{../../others/#3}
}}

\newcounter{mypage}
\newcommand{\fromBook}[6][scale=0.5]{
\setcounter{mypage}{20}
\addtocounter{mypage}{#2}
\begin{tabular}{@{}r@{}}
\includegraphics[page=\themypage, trim = #3 #4 #5 #6, clip, #1]%
  {../../learning-concurrent-scala/learning-concurrent-programming-in-scala.pdf}
  \\[0mm]
  \textcolor{myGray}{$\left[\begin{array}{@{}r@{}}
  ~\\[-7mm]
  \text{{\tiny in \emph{``Learning Concurrent}}}\\[-3mm]
  \text{{\tiny \emph{Programming in Scala''}, pg.\,#2}}
  \\[0mm]
  \end{array}\right]$}
\end{tabular}
}

\newcommand{\fromBookW}[4][width=\textwidth]{
\setcounter{mypage}{20}
\addtocounter{mypage}{#2}
\begin{tabular}{@{}r@{}}
\includegraphics[page=\themypage, trim = 20mm #3 20mm #4, clip, #1]%
  {../../learning-concurrent-scala/learning-concurrent-programming-in-scala.pdf}
  \\[-2mm]
  \textcolor{myGray}{$\left[\begin{array}{@{}r@{}}
  ~\\[-7mm]
  \text{{\tiny in \emph{``Learning Concurrent}
                      \emph{Programming in Scala''}, pg.\,#2}}
  \\[0mm]
  \end{array}\right]$}
\end{tabular}
}

\newcommand{\fromBookWframed}[4][width=\textwidth]{
\setcounter{mypage}{20}
\addtocounter{mypage}{#2}
\begin{tabular}{@{}r@{}}
\includegraphicsframed[page=\themypage, trim = 20mm #3 20mm #4, clip, #1]%
  {../../learning-concurrent-scala/learning-concurrent-programming-in-scala.pdf}
  \\[-2mm]
  \textcolor{myGray}{$\left[\begin{array}{@{}r@{}}
  ~\\[-7mm]
  \text{{\tiny in \emph{``Learning Concurrent}
                      \emph{Programming in Scala''}, pg.\,#2}}
  \\[0mm]
  \end{array}\right]$}
\end{tabular}
}

\newcommand{\fromAuthor}[3][width=\textwidth]{
\begin{tabular}{@{}r@{}}
\includegraphics[#1]%
  {src/img/#3}
  \\[-2mm]
  \textcolor{myGray}{$\left[\begin{array}{@{}r@{}}
  ~\\[-7mm]
  \text{{\tiny in #2}}
  \\[0mm]
  \end{array}\right]$}
\end{tabular}
}


%-------------- template --------------------------------------------------
 %!TEX root=../../1-intro.tex
 \usetheme{metropolis}
\usepackage{appendixnumberbeamer}

%------ Setting lecture info ----------------------------------------------
\newcounter{lectureID}
\stepcounter{lectureID}
\newcommand{\getLecture}{\arabic{lectureID}\xspace}
\newcommand{\setLectureBasic}[1]{
  \title{
    #1
    }
  \author{Nelma Moreira ~~\&~~ \textbf{Jos\'{e} Proen\c{c}a}}
  \institute{CISTER -- U.Porto, Porto, Portugal
            \hfill 
            \begin{tabular}{r@{}}
            \url{https://fm-dcc.github.io/cp2425}
            \end{tabular}
            }
  \date{Concurrent programming (CC3040) 2024/2025}
  % logos of institutions
  \titlegraphic{
    \begin{textblock*}{5cm}(2.0cm,7.00cm)
       \includegraphics[scale=0.14]{src/img/logos/fcup}\hspace*{.85cm}~%
    \end{textblock*}
    \begin{textblock*}{5cm}(6.0cm,7.25cm)
      \includegraphics[scale=0.43]{src/img/logos/cmup}
    \end{textblock*}
    \begin{textblock*}{5cm}(10.4cm,7.45cm)
      \includegraphics[scale=0.16]{src/img/logos/cister}
    \end{textblock*}
  }  
  \frame[plain]{\titlepage}
}
\newcommand{\setLecture}[2]{\setcounter{lectureID}{#1}\setLectureBasic{#1. #2}}

%------ Counters for exercises ----------------------------------------------
\newcounter{cExercise}
\newcommand{\exercise}{\stepcounter{cExercise}Ex.\,\arabic{lectureID}.\arabic{cExercise}:\xspace}
\newcommand{\exerciseBack}{\addtocounter{cExercise}{-1}}
\newcommand{\exerciseAdd}{\stepcounter{cExercise}}
\newcommand{\doExercise}[3][0mm]{\begin{exampleblock}{\exercise #2}\wrap{\rule{0pt}{#1}}#3\end{exampleblock}}
\newcommand{\doSimpleExercise}[2][0mm]{\begin{exampleblock}{}\wrap{\rule{0pt}{#1}}\structure{\textbf{\exercise} #2}\end{exampleblock}}

% Slide
\newenvironment{slide}[1]{\begin{frame}\frametitle{#1}}{\end{frame}}



% Base colors (from metropolis theme)
\definecolor{metDarkBrown}{HTML}{604c38}
\definecolor{metDarkTeal}{HTML}{23373b}
\definecolor{metLightBrown}{HTML}{EB811B}
\definecolor{metLightGreen}{HTML}{14B03D}
\definecolor{metDarkGreen}{HTML}{4f9277}

 

\metroset{numbering=fraction,progressbar=frametitle}

% \setbeamercolor*{structure}{fg=blue!80!black}
\setbeamercolor*{structure}{fg=metLightGreen}

% % \definecolor{MainColour}{rgb}{0., 0.25, 0.8}
% \colorlet{MainColour}{blue!50!black}
% \colorlet{BgColour}{blue!10}
% \colorlet{BarColour}{blue!50!black}

% %\usetheme{CambridgeUS}%{Copenhagen}%{Frankfurt}%{Singapore}%{CambridgeUS}
% \usecolortheme[named=MainColour]{structure} 
% \useoutertheme[subsection=false]{miniframes}
% \useinnertheme{circles}
% %\useinnertheme[shadow=false]{rounded}
% \setbeamertemplate{blocks}[rounded][shadow=false]

% \setbeamercovered{transparent} 
% \setbeamertemplate{navigation symbols}{} %Remove navigation bar
% \setbeamertemplate{footline}[frame number] % add page number
% \setbeamercolor{postit}{fg=MainColour,bg=BgColour}
% \setbeamercolor{structure}{bg=black!10}
% %\setbeamercolor{palette primary}{use=structure,fg=red,bg=green}
% %\setbeamercolor{palette secondary}{use=structure,fg=red!75!black,bg=green}
% \setbeamercolor{palette tertiary}{use=structure,bg=BarColour,fg=white}
% %\setbeamercolor{palette quaternary}{fg=black,bg=green}
% %\setbeamercolor{normal text}{fg=black,bg=white}
% %\setbeamercolor{block title alerted}{fg=red,bg=green}
% %\setbeamercolor{block title example}{bg=black!10,fg=green}
\setbeamercolor{block body}{bg=black!5}

% \setbeamercolor{block title alerted}{bg=red!25}
% \setbeamercolor{block body alerted}{bg=red!10}

% \setbeamercolor{block title example}{bg={rgb:green,2;black,1;white,5}}
\setbeamercolor{block body example}{bg={rgb:green,2;black,1;white,20}}
\setbeamercolor{block body alerted}{bg={metLightBrown!25}}

% \setbeamertemplate{itemize item}{\color{black!10}$\blacksquare$}
\setbeamercolor{itemize item}{fg=metDarkTeal}
\setbeamercolor{itemize subitem}{fg=metDarkTeal}

\setbeamercolor{graybc}{fg=black,bg=black!10}
\newcommand{\myblock}[1]{\begin{beamercolorbox}[dp=1ex,center,rounded=true]%
  {graybc} {\large \textbf{#1}} \end{beamercolorbox}}%

%%%%%%%%%

% \definecolor{barcolor}{rgb}{.65,.79,.92} % FCUP color
\colorlet{barcolor}{metDarkGreen!60}


% Configuring the foot line
\setbeamercolor{author in head/foot}{fg=metDarkTeal, bg=barcolor}%
\setbeamercolor{date in head/foot}{fg=barcolor!75!black}%
\setbeamertemplate{footline}
{
  \leavevmode%
  \hbox{%
  \begin{beamercolorbox}[wd=.4\paperwidth,ht=2.25ex,dp=1ex,center]{author in head/foot}%
    \usebeamerfont{author in head/foot}\insertshortauthor
  \end{beamercolorbox}%
  \begin{beamercolorbox}[wd=.5\paperwidth,ht=2.25ex,dp=1ex,center]{title in head/foot}%
    \usebeamerfont{title in head/foot}\insertsection
  \end{beamercolorbox}%
  \begin{beamercolorbox}[wd=.1\paperwidth,ht=2.25ex,dp=1ex,right]{date in head/foot}%
    \insertframenumber{} / \inserttotalframenumber\hspace*{2ex} 
  \end{beamercolorbox}}%
  \vskip0pt%
}
% No configuration symbols
\setbeamertemplate{navigation symbols}{}


%%%%%%% Custom bar above

\setbeamercolor{frametitle}{bg=barcolor,fg=metDarkTeal}
\setbeamercolor{progress bar}{fg=barcolor!75!black}

%%% Custom bar with a LOGO

\makeatletter
\setbeamertemplate{frametitle}{%
  \nointerlineskip%
  \begin{beamercolorbox}[%
      wd=\paperwidth,%
      sep=0pt,%
      leftskip=\metropolis@frametitle@padding,%
      rightskip=\metropolis@frametitle@padding,%
    ]{frametitle}%
  \metropolis@frametitlestrut@start%
  \insertframetitle%
  \nolinebreak%
  \metropolis@frametitlestrut@end%
  \hfill
  \raisebox{-1.3ex}{\includegraphics[height=4ex,keepaspectratio]{src/img/logos/fcup-bar}}
  \end{beamercolorbox}%
}
\makeatother





%----------------------------------------------------------------------------

\begin{document}

\setLecture{5}{Actor model using the Akka framework}


\section{Overview}

\begin{frame}{We are here}

  \vspace*{-2mm}

  \begin{block}{Blocks of sequential code running concurrently and sharing memory:}
    
  \begin{itemize}
    \item What is Scala?
    \item Concurrency in Java and its memory model
    \item Basic concurrency blocks and libraries
    % \item Futures and promises
    \item \textcolor{gray}{\emph{Futures and Promises}}
    \item \textcolor{gray}{\emph{Data-Parallel Collections}}
    \item \textcolor{gray}{\emph{Reactive Programming (Concurrently)}}
    \item \textcolor{gray}{\emph{Software Transactional Memory}}
    \alert{\item Actor model}
  \end{itemize}
  \end{block}
\end{frame}



% \fromBookW[scale=0.7]{32}{98mm}{39mm}
% \fromBook[scale=0.65]{32}{43mm}{98mm}{43mm}{39mm}

% \begin{frame}[fragile]\frametitle{Current thread}
% ~\\[-8mm]
% \begin{columns}
% \begin{column}{0.49\textwidth}
% \begin{lstlisting}
% ...
% \end{lstlisting}
% \end{column}
% \begin{column}{0.49\textwidth}
% ...
% \end{column}
% \end{columns}
% \end{frame}


\begin{frame}\frametitle{What is the actor model}
  \splittwo{0.4}{0.6}{
    \begin{itemize}
      \item \alert{Asynchronous} message exchange between actors
      \item Introduced in \alert{Erlang}\\(we use Akka's actor library)\\~
      \item<2> Active, autonomous, \alert{no shared memory}, no synchronisation
    \end{itemize}
  }{
    \centering
    \includegraphicsframed[scale=1]{src/img/actor.pdf}\\~\\
    \only<2>{\includegraphicsframed[scale=1]{src/img/actor-com.pdf}}
  }
\end{frame}

\begin{frame}[t]\frametitle{What we will see}

  We will use the Akka framework for actors for:
  \begin{itemize}
    \item Declaring \alert{actor classes} and creating \alert{actor instances}
    \item Modelling \alert{actor state} and complex \alert{actor behaviours}
    \item Manipulating the \alert{actor hierarchy} and the \alert{actor lifecycle}
    \item The different message-passing patterns used in \alert{actor communication}
    \item Error recovery using the built-in \alert{actor supervision} mechanism
    \item Using \alert{remote actors} to build concurrent and distributed programs
  \end{itemize}

  Documentation: \url{https://doc.akka.io/docs/akka}


\end{frame}


\section{Creating actors}

\begin{frame}[t]\frametitle{Core concepts}
  ~\\[-8mm]    
  \splittwo{0.5}{0.5}{
    \begin{exampleblock}{Actor system}
      Hierarchical group of actors with shared configurations, supporting actor creation and logging.
    \end{exampleblock}

    \begin{exampleblock}{Actor class}
      Template that describes the states and behaviour of an actor, used to create instances.
    \end{exampleblock}

    \begin{exampleblock}{Actor instance}
      Entity that exists at runtime, with a state and capable of sending and receiving messages.
    \end{exampleblock}

    \begin{exampleblock}{Message}
      Unit of communication that actors use to communicate. In Akka, any object can be a message.
    \end{exampleblock}
  }{
    \begin{exampleblock}{Mailbox}
      Memory block that is used to buffer messages for a given actor instance.
    \end{exampleblock}

    \begin{exampleblock}{Actor reference}
      Object that allows an object to send messages to a specific actor instance.
    \end{exampleblock}

    \begin{exampleblock}{Dispatcher}
      Component that decides when actors are allowed to process messages. In Akka every dispatcher is also an execution context.
    \end{exampleblock}
  }
\end{frame}

\begin{frame}[fragile]\frametitle{My first actor (class) in Akka}
~\\[-8mm]
\begin{columns}
\begin{column}{0.56\textwidth}
\begin{lstlisting}[emph={Actor, Logging, receive, system,stop}]
import akka.actor._
import akka.event.Logging

class HelloActor(val hello: String) extends Actor {
  val log = Logging(context.system, this)
  def receive = { // Any => Unit (partial)
    case `hello` =>
      log.info(
        s"Received a '$hello'... $hello!")
    case msg     =>
      log.info(
        s"Unexpected message '$msg'")
      context.stop(self)
  }
}
\end{lstlisting}
\end{column}
\begin{column}{0.44\textwidth}
\begin{itemize}
  \item Each \cod{HelloActor} receives messages
  \item ... if it receives its \cod{hello}, it logs and \structure{continues}
  \item ... if it receives something else, it \structure{stops}
  \item \cod{context} -- provides core functions, such as \cod{stop}
  \item \cod{self} -- is the instance's actor reference
\end{itemize}
\end{column}
\end{columns}
\end{frame}


\begin{frame}[fragile]\frametitle{Configuring an actor in Akka}
~\\[-6mm]
\begin{columns}
\begin{column}{0.51\textwidth}
\begin{lstlisting}[emph={Actor, Logging, Props}]
object HelloActor { // companion
  // two factory methods below
  def props(hello: String) =
    Props(new HelloActor(hello))
  def propsAlt(hello: String) =
    Props(classOf[HelloActor], hello)
  //def propsAlt2 = Props[HelloActor]
}
\end{lstlisting}
\begin{exampleblock}{Actor configuration}
  \begin{itemize}
      \item actor class
      \item constructor arguments
      \item mailbox
      \item dispatcher
  \end{itemize} 
\end{exampleblock}
\end{column}
\begin{column}{0.49\textwidth}
\begin{block}{Props}
  \begin{itemize}
    \item can receive a block of code,
     used each time a new actor instance is created;
     \item can receive a \cod{Class} object and its arguments
     \item can be sent over the network (should be self-contained)
     \item avoid creating \cod{Props} in the actor class, and use factory methods instead
  \end{itemize}
\end{block}
\end{column}
\end{columns}
\end{frame}

\begin{frame}[fragile]\frametitle{My first actor system with an instance}
\begin{lstlisting}[emph={Actor, Logging, Props, ourSystem}]
// in build.sbt:
libraryDependencies ++= Seq( ...
  ,"com.typesafe.akka" %% "akka-actor" % "2.8.5"
  ,"com.typesafe.akka" %% "akka-remote" % "2.8.5"
)
\end{lstlisting}
\begin{lstlisting}[emph={Actor, Logging, Props, ourSystem}]
lazy val ourSystem = akka.actor.ActorSystem("OurExampleSystem")
\end{lstlisting}
\begin{lstlisting}[emph={Actor, Logging, Props, ourSystem,terminate,sleep,actorOf}]
object ActorsCreate extends App {
  val hiActor: ActorRef =
    ourSystem.actorOf(HelloActor.props("ola"), name = "greeter")
  hiActor ! "ola"
  Thread.sleep(1000)
  hiActor ! "hi"
  Thread.sleep(1000)
  ourSystem.terminate()
}
\end{lstlisting}
\end{frame}


\begin{frame}[fragile]\frametitle{HelloActor in one slide}
~\\[-6mm]
\begin{columns}
\begin{column}{0.58\textwidth}
\begin{lstlisting}[emph={Actor, Logging, Props, receive, stop, log, terminate, actorOf}]
import akka.actor._
import akka.event.Logging

object ActorsCreate2 extends App {
  lazy val ourSystem =
    akka.actor.ActorSystem("OurSystem")

  class Hi(val hi: String) extends Actor {
    val log = Logging(context.system, this)
    def receive = {
      case `hi` =>
        log.info(s"Got a '$hi'... $hi!")
      case msg     =>
        log.info(s"Unexpected '$msg'")
        context.stop(self)
    }
  }
\end{lstlisting}
\end{column}
\begin{column}{0.42\textwidth}
\begin{lstlisting}[emph={Actor, Logging, Props, receive, stop, log, terminate, actorOf},mathescape=true]
  val hiActor: ActorRef =
    ourSystem.actorOf(
      Props(new Hi("ola")),
      name = "greeter")

  hiActor ! "ola"
  // ... Got a 'ola'... ola!
  Thread.sleep(1000)
  hiActor ! "yo"
  // ... Unexpected 'yo'
  Thread.sleep(1000)
  hiActor ! "$\textsf{\textcolor{green!60!black}{\foreignlanguage{russian}{привет}}}$" 
  // ... Message (...) was not delivered.
  Thread.sleep(1000)
  ourSystem.terminate()
}
\end{lstlisting}
\end{column}
\end{columns}
\end{frame}

\begin{frame}[fragile]\frametitle{HelloActor's feedback}
~\\[-6mm]
\begin{lstlisting}[emph={Actor, Logging, Props, receive, stop, log, terminate, actorOf},mathescape=true]
  hiActor ! "ola"
  // ... [akka://OurSystem/user/greeter] Got a 'ola'... ola!

  Thread.sleep(1000)
  hiActor ! "yo"
  // ... [akka://OurSystem/user/greeter] Unexpected 'yo'

  Thread.sleep(1000)
  hiActor ! "$\textsf{\textcolor{green!60!black}{\foreignlanguage{russian}{привет}}}$" 
  // ... [akka://OurSystem/user/greeter] Message [...] to Actor[akka://OurSystem/user/greeter#-726408098] was not delivered.

  Thread.sleep(1000)
  ourSystem.terminate()
}
\end{lstlisting}
\end{frame}


\begin{frame}[fragile]\frametitle{Unhandled messages?}

\begin{lstlisting}[emph={Actor, Logging, Props, ourSystem,terminate,sleep,actorOf,receive,unhandled}]
class DeafActor extends Actor {
  val log = Logging(context.system, this)
  def receive = PartialFunction.empty
  // default: ignore and log
  override def unhandled(msg: Any) = msg match {
    case msg: String => log.info(s"I do not hear '$msg'")
    case msg         => super.unhandled(msg)
  }
}
\end{lstlisting}
    
\begin{lstlisting}[emph={Actor, Logging, Props, ourSystem,terminate,sleep,actorOf}]
object ActorsUnhandled extends App {
  val deafActor: ActorRef =
    ourSystem.actorOf(Props[DeafActor], name = "deafy")
  deafActor ! "ola"
  Thread.sleep(1000)
  deafActor ! 1234
  Thread.sleep(1000)
  ourSystem.terminate()
}
\end{lstlisting}
\end{frame}


\section{Modelling actor behaviour}

\begin{frame}[fragile]\frametitle{My 2nd example in Akka: a (stateful) countdown}
~\\[-6mm]
\begin{columns}
\begin{column}{0.48\textwidth}
\only<2>{\alert{Not allowed in Akka:}}
\begin{lstlisting}[emph={Actor, Logging, Props, ourSystem,terminate,sleep,actorOf,receive}]
class CountdownActor extends Actor{
  var n = 10
            // never do this
  def receive = if (n > 0) {
    case "count" =>
      log(s"n = $n")
      n -= 1
  } else PartialFunction.empty
}
\end{lstlisting}
\end{column}
\begin{column}{0.52\textwidth}
\pause
\alert{Correct in Akka, using \cod{become}:}
\begin{lstlisting}[emph={Actor, Logging, Props, ourSystem,terminate,sleep,actorOf,receive,become}]
class CountdownActor extends Actor {
  val log = Logging(context.system, this)
  var n = 10
  def counting: Actor.Receive = {
    case "count" =>
      n -= 1
      log.info(s"n = $n")
      if (n == 0) context.become(done)
  }
  def done = PartialFunction.empty
  def receive = counting
}
\end{lstlisting}
\end{column}
\end{columns}
\end{frame}

\begin{frame}\frametitle{Actor as a transition system}
  \centering
  \fromBookW[scale=0.7]{278}{126mm}{43mm}
\end{frame}

\begin{frame}[fragile]\frametitle{Running the countdown}

\begin{lstlisting}[emph={Actor, Logging, Props, ourSystem,terminate,sleep,actorOf,receive,become}]
object ActorsCountdown extends App {
  val countdown = ourSystem.actorOf(Props[CountdownActor])
  for (i <- 0 until 20) countdown ! "count"
  Thread.sleep(1000)
  ourSystem.terminate()
}
\end{lstlisting}


\end{frame}


\section{Actor hierarchy and lifecycle}

\begin{frame}[fragile]\frametitle{New example with a parent}

\begin{columns}
\begin{column}{0.47\textwidth}
~\\[1mm]
\centering
\begin{tikzpicture}
  \tikzstyle{nd}=[rectangle,inner sep=4pt,minimum width=10mm,draw=black!60,fill=barcolor!20,very thick]
  \node[nd,draw=none,fill=none,inner sep=1pt](sys){``OurExampleSystem''};
  \node[nd, below=2mm of sys](par){``parent''};
  \node[nd, below=2mm of par,xshift=-12mm](a){child$_1$};
  \node[nd, below=2mm of par,xshift= 12mm](b){child$_2$};
  \draw[<-,thick] (sys)edge(par) (par)edge(a)edge(b); 
\end{tikzpicture}

\pause
\begin{lstlisting}[emph={Actor, Logging, Props, ourSystem,terminate,context,actorOf,receive}]
class ChildActor extends Actor {
  val log = Logging(context.system, this)
  def receive = {
    case "sayhi" =>
      val parent = context.parent
      log.info(s"my parent $parent made me say hi!")
  }
  override def postStop() {
    log.info("child stopped!")
  }
}
\end{lstlisting}
\end{column}
\begin{column}{0.54\textwidth}
\begin{lstlisting}[emph={Actor, Logging, Props, ourSystem,terminate,sleep,actorOf,receive,become,context}]
class ParentActor extends Actor {
  val log = Logging(context.system, this)
  def receive = {
    case "create" =>
      context.actorOf(Props[ChildActor])
      log.info(s"created a kid; children = ${context.children}")
    case "sayhi" =>
      log.info("Kids, say hi!")
      for (c <- context.children)
        c ! "sayhi"
    case "stop" =>
      log.info("parent stopping")
      context.stop(self)
  }
}
\end{lstlisting}
\end{column}
\end{columns}
\end{frame}


\begin{frame}[fragile]\frametitle{A more complete view of the hierarchy}
~\\[-6mm]
\begin{columns}
\begin{column}{0.57\textwidth}
\begin{lstlisting}[emph={Actor, Logging, Props, ourSystem,terminate,sleep,actorOf,receive}]
object ActorsHierarchy extends App {
  val parent = ourSystem.actorOf(Props[ParentActor], "parent")
  parent ! "create"
  parent ! "create"
  Thread.sleep(1000)
  parent ! "sayhi"
  Thread.sleep(1000)
  parent ! "stop"
  Thread.sleep(1000)
  ourSystem.terminate()
}\end{lstlisting}
\end{column}
\begin{column}{0.41\textwidth}
  \fromBook[scale=0.6]{284}{46mm}{30mm}{46mm}{143mm}
  ~\\
  \splittwo{0.5}{0.5}{
      - ActorSystem
    \\- sys.terminate
    \\- sys/ctxt.actorOf
    \\
  }{
      - ctxt.stop
    \\- ctxt.become
    \\\alert{- ctxt.children}
    \\\alert{- ctxt.parent}
  }
\end{column}
\end{columns}
\end{frame}


\begin{frame}[fragile]\frametitle{A more complete view of the hierarchy}
~\\[-6mm]
\begin{columns}
\begin{column}{0.57\textwidth}
\begin{itemize}
  \item \alert{parent} actor stops $\Rightarrow$ its \alert{children} stop
  \item \structure{user} and \structure{system}:
    \\ are \alert{guardian actors} -- at the top of the hierarchy, to log, restart actors, etc.
  \item hierarchy visible when printing an actor ref, e.g., for the first child;
    \\\code{akka://OurExampleSystem/user/parent/}\texttt{\$}\code{a}
  \item \pause \alert{Next:}
    \cod{ctxt.actorSelection(path)}
\end{itemize}
\end{column}
\begin{column}{0.41\textwidth}
  \fromBook[scale=0.6]{284}{46mm}{30mm}{46mm}{143mm}
  ~\\
  \splittwo{0.5}{0.5}{
      - ActorSystem
    \\- sys.terminate
    \\- sys/ctxt.actorOf
    \\
  }{
      - ctxt.stop
    \\- ctxt.become
    \\- ctxt.children
    \\- ctxt.parent
  }
\end{column}
\end{columns}
\end{frame}


\begin{frame}[fragile]\frametitle{Discovering actors in the hierarchy}

\begin{lstlisting}[emph={Actor, Logging, Props, ourSystem,terminate,sleep,actorOf,receive,Identify,actorSelection}]
class CheckActor extends Actor {
  val log = Logging(context.system, this)
  def receive = {
    case path: String =>
      log.info(s"checking path $path")
      context.actorSelection(path) ! Identify(path)
    case ActorIdentity(path, Some(ref)) =>
      log.info(s"found actor $ref at $path")
    case ActorIdentity(path, None) =>
      log.info(s"could not find an actor at $path")
  }
}
\end{lstlisting}    

\pause
\begin{lstlisting}[emph={Actor, Logging, Props, ourSystem,terminate,sleep,actorOf,receive,Identify,actorSelection}]
val checker = ourSystem.actorOf(Props[CheckActor], "checker")

checker ! "../*"      // finds the checker and its siblings
checker ! "../../*"   // finds user and system guardians
checker ! "/system/*" // finds internal actors
checker ! "/user/checker2" // logs that no actors were found
\end{lstlisting}
\end{frame}


\begin{frame}\frametitle{Once an actor throws an exception...}
  When an actor throws an exception, a new \alert{``replacement'' actor} is created,
  with the same:
  \begin{itemize}
    \item arguments
    \item mailbox
    \item ActorRef
    \pause
      \begin{itemize}
        \item hence never leak the actual \alert{this} reference!
      \end{itemize}
  \end{itemize}
\end{frame}

\begin{frame}\frametitle{Actor lifecycle}

\begin{itemize}
  \item Created -- \cod{actorOf}
  \item Before starting to process messages -- \cod{preStart()}
  \item After an exception -- \cod{preRestart(t: Throwable, msg: Option[Any])}
    \begin{itemize}
      \item before creating a new actor
      \item when all children are stopped
    \end{itemize}
  \item After recreating a restarted actor -- \cod{postRestart(t: Throwable)}
    \begin{itemize}
      \item the new actor is then assigned the previous mailbox
    \end{itemize}
  \item After an actor terminates -- \cod{postStop()}
    \begin{itemize}
      \item called by the default implementation of \code{preRestart}
    \end{itemize}
\end{itemize}
\end{frame}

\begin{frame}\frametitle{Actor lifecycle in a diagram}
  \centering
  \fromBookW[scale=0.9]{289}{35mm}{130mm}
\end{frame}




\section{Synchrony vs. Asynchrony}

\begin{frame}[t]\frametitle{Sending x and y from A to B}
~\\[-8mm]
\splittwo{0.5}{0.5}{
\begin{exampleblock}{Synchronous (as in CCS)}
$A = x!\,.\,y!$\\
$B = x?\,.\,y?$\\
$A~|~B\backslash\{x,y\}$
\\[4mm]
\pause
$\Rightarrow ~~~~ \tau_x . \tau_y$
\pause
\end{exampleblock}
}{
\begin{alertblock}{Asynchronous (as in Akka)}
$x!$ happens before $y!$\\
$x?$ happens before $y?$\\\pause
$x!$ happens before $x?$\\
$y!$ happens before $y?$\\\pause
$y!$ ?? $x?$
\end{alertblock}
}
~\\
\splittwo{0.5}{0.5}{
\pause
Different formalisations for global beh.:
\begin{itemize}
\item Message sequence charts
\item Event structures
\item Automata over interactions
\item Choreographies:
  \\$A\to B:x~~;~~A\to B:y$
\end{itemize}
}{
  No duplication
  \\ No messages lost
  \\ No messages reordered
  \\ No blocking send
  \\[4mm]
  \structure{Synchrony} modelled with \alert{Asynchrony}?
  \\and vice-versa?
}
\end{frame}

\begin{frame}[t]\frametitle{Diamond problem: sending by two routes}
    
    $A\to B: x;$\\
    $A\to C: y;$\\
    $C\to B: z$

~\\

$B$ must be ready to receive `$x?$' and `$z?$' by any order

\end{frame}



\section{Error recovery with actors}

\begin{frame}[t]\frametitle{Stopping an actor}
\begin{alertblock}{Main ways to stop an actor:}  
\begin{itemize}
  \item \cod{context.stop(act)} -- \alert{stops} \cod{act} and its children, once it finishes processing their current message
  \item \cod{Kill} message -- \alert{restarts} the target actor once it is received
    %once received, will restart the target actor without losing its mailbox
  \item \cod{PoisonPill} message -- \alert{stops} the target actor \structure{after} once it is processed
\end{itemize}
\end{alertblock}

\pause

\begin{block}{Stopping in more complex scenarios:}
  \begin{itemize}
    \item Using Akka's \alert{DeathWatch} (next slide)
  \end{itemize}
\end{block}
\end{frame}

\begin{frame}[t]\frametitle{Pingy-Pongy example}
% // pongy: pingy?"ping";  pingy!"pong"; stop(self)
% // pingy: master?pongy; pongy!"ping"; pongy?reply; master!reply; loop
% // master: sys?"start"; pingy!pongy; _?"pong"; stop
% master->pingy:pongyRef;
% pingy->pongy:ping;
% pongy->pingy:pong;
% pongy:stop;
% pingy->master:pong;
% master:stop
\centering    
\begin{tabular}{cc}
  \wrap{\includegraphicsframed[width=0.5\textwidth]{src/img/pingy-pongy.pdf}}
  &
  \begin{minipage}{0.3\textwidth}
  \begin{itemize}
    \item Example used in the book to illustrate the ask-reply pattern
    \item (in pingy: \code{val reply = pongy ? "ping"})
    \item We will adapt it for a graceful shutdown
  \end{itemize}
  \end{minipage}
\end{tabular}
\end{frame}

\begin{frame}[fragile]\frametitle{Graceful Pingy-Pongy}
~\\[-6mm]
\begin{columns}[t]
\begin{column}{0.56\textwidth}
% \only<2>{\alert{Not allowed in Akka:}}
\begin{lstlisting}[emph={Actor, Logging, Props, ourSystem,terminate,sleep,actorOf,receive, watch,stop,Terminated}]
class GracefulPingy extends Actor {
  val log = Logging(context.system, this)
  val pongy =
    context.actorOf(Props[Pongy], "pongy")
  context.watch(pongy)
  def receive = {
    case "start" => pongy ! "ping"
    case "pong"  => log.info("Got a pong")
    case "Die, Pingy!" =>
      context.stop(pongy)
    case Terminated(`pongy`) =>
      context.stop(self)
} }
\end{lstlisting}
\end{column}
\begin{column}{0.44\textwidth}
\begin{lstlisting}[emph={Actor, Logging, Props, ourSystem,terminate,sleep,actorOf,receive,become,stop}]
class Pongy extends Actor {
  val log =
    Logging(context.system,this)
  def receive = {
    case "ping" =>
      log.info("Got a ping -- ponging back!")
      sender ! "pong"
  }
  override def postStop() = log.info("pongy going down")
}
\end{lstlisting}
\end{column}
\end{columns}
\end{frame}

\begin{frame}[fragile]\frametitle{Running the gracefull app}
\centering
\splittwo{0.4}{0.6}{
  \includegraphicsframed[width=\textwidth]{src/img/gpingy-pongy.pdf}
}{

  \begin{exampleblock}{Mechanism 1 (pingy $\leftrightarrow$ pongy)}
  \begin{itemize}
    \item \cod{context.watch(pongy)} -- the \code{DeathWatch}
    \item wait for \cod{Terminated} message
  \end{itemize}
  \end{exampleblock}
  \begin{alertblock}{Mechanism 2 (ourSystem $\leftrightarrow$ pingy)}
  \begin{itemize}
    \item ask to \code{``Die''}
    \item check if it terminated -- using Futures
  \end{itemize}
  \end{alertblock}
}

\end{frame}

\begin{frame}[fragile]\frametitle{Running the gracefull app (code)}
~\\[-6mm]
% \begin{columns}[t]
% \begin{column}{0.56\textwidth}
% \begin{lstlisting}[emph={Actor, Logging, Props, ourSystem,terminate,sleep,actorOf,receive, watch,stop}]
%   val log = Logging(context.system, this)
%   val pongy =
%     context.actorOf(Props[Pongy], "pongy")
%   context.watch(pongy)
%   def receive = {
%     case "start" => pongy ! "ping"
%     case "pong"  => log.info("Got a pong")
%     case "Die, Pingy!" =>
%       context.stop(pongy)
%     case Terminated(`pongy`) =>
%       context.stop(self)
% } }
% \end{lstlisting}
% \end{column}
% \begin{column}{0.44\textwidth}
\begin{lstlisting}[emph={Actor, Logging, Props, ourSystem,terminate,sleep,actorOf,receive,become,stop,onComplete,gracefulStop}]
import akka.pattern.gracefulStop

object CommunicatingGracefulStop extends App {
  val gracePingy = ourSystem.actorOf(Props[GracefulPingy], "gracePingy")
  gracePingy ! "start"

  val stopped = gracefulStop(gracePingy, 3.seconds, "Die, Pingy!")
  stopped onComplete { // stopped is a Future (not covered)
    case Success(x) =>
      log("graceful shutdown successful")
      ourSystem.terminate()
    case Failure(t) =>
      log("grace not stopped!")
      ourSystem.terminate()
} }
\end{lstlisting}
% \end{column}
% \end{columns}
\end{frame}


\begin{frame}[fragile]\frametitle{Handling children's exceptions (Actor supervision)}
~\\[-10mm]
\begin{columns}[t]
\begin{column}{0.54\textwidth}
\begin{lstlisting}[emph={Actor, Logging, Props, ourSystem,terminate,sleep,actorOf,receive, watch,stop,Terminated}]

class Naughty extends Actor {
  val log = Logging(context.system,this)
  def receive = {
    case s: String => log.info(s)
    case msg => throw new RuntimeException
  }
  override def postRestart(t:Throwable)=
    log.info("naughty restarted")
}
\end{lstlisting}
\end{column}
\begin{column}{0.50\textwidth}
\begin{lstlisting}[emph={Actor, Logging, Props, ourSystem,terminate,sleep,actorOf,receive,become,stop,supervisorStrategy}]
import SupervisorStrategy._
class Supervisor extends Actor {
  context.actorOf(Props[Naughty], "naughty")
  def receive = PartialFunction.empty
  override val supervisorStrategy =
    OneForOneStrategy() {
      case ake: ActorKilledException => Restart
      case _ => Escalate
} }
\end{lstlisting}
\end{column}
\end{columns}
~\\[-4pt]
\begin{lstlisting}[emph={Actor, Logging, Props, ourSystem,terminate,sleep,actorOf,receive,become,stop}]
ourSystem.actorOf(Props[Supervisor], "super")
val children = ourSystem.actorSelection("/user/super/*")
children ! "hello" // succeeds
children ! Kill    // stops naughty, but super restarts it
children ! "sorry about that" // succeeds
children ! "kaboom".toList    // naughty and super throw exception
\end{lstlisting}
\end{frame}


\section{Remote actors over TCP}

\begin{frame}[fragile]\frametitle{Compilation with remote actors}
~\\[-6.5mm]
\begin{tabular}{ll}
\begin{minipage}{0.18\textwidth}
 \cod{build.sbt} needs to import akka-remote:
\end{minipage}
&
\begin{minipage}{0.82\textwidth}
\begin{lstlisting}[emph={Actor, Logging, Props, ourSystem,terminate,sleep,actorOf,receive,become,stop}]
libraryDependencies ++= Seq(
   ...
  ,"com.typesafe.akka" %% "akka-actor" % "2.8.5" // or older
  ,"com.typesafe.akka" %% "akka-remote" % "2.8.5"
)
\end{lstlisting}
\end{minipage}
\\
\begin{minipage}{0.18\textwidth}
  Network configured with Netty library
\end{minipage}
&
\begin{minipage}{0.82\textwidth}
\begin{lstlisting}[
  showstringspaces=false,
  emph={Actor, Logging, Props, ourSystem,terminate,sleep,actorOf,receive,become,stop,remotingConfig,remotingSystem}]
import com.typesafe.config._
def remotingConfig(port: Int) = ConfigFactory.parseString(s"""
  akka {
    actor.provider = "akka.remote.RemoteActorRefProvider"
    remote {
      enabled-transports = ["akka.remote.netty.tcp"]
      netty.tcp {
        hostname = "127.0.0.1"
        port = $port }
    }
  }""")
def remotingSystem(name: String, port: Int): ActorSystem =
     ActorSystem(name, remotingConfig(port))
\end{lstlisting}
\end{minipage}
% \end{itemize}
\end{tabular}
\end{frame}


\begin{frame}[fragile]\frametitle{Remote Pingy-Pongy -- running two Apps!}
~\\[-6mm]
\begin{columns}[t]
\begin{column}{0.46\textwidth}
\begin{lstlisting}[basicstyle=\ttfamily\tiny,emph={Actor, pingy, pongy, runner, Logging, Props, ourSystem,terminate,sleep,actorOf,receive, watch,stop,Terminated}]
object RemotingPongySystem extends App {
  val system = remotingSystem("PongyDimension", 24321)
  val pongy = system.actorOf(Props[Pongy], "pongy")
  Thread.sleep(15000)
  system.terminate()
}
\end{lstlisting}
\begin{lstlisting}[basicstyle=\ttfamily\tiny,
  emph={Actor, pingy, pongy, runner, Logging, Props, ourSystem,terminate,sleep,actorOf,receive, watch,stop,Terminated}]
object RemotingPingySystem extends App {
  val system = remotingSystem("PingyDimension", 24567)
  val runner = system.actorOf(Props[Runner], "runner")
  runner ! "start"
  Thread.sleep(5000)
  system.terminate()
}
\end{lstlisting}
\end{column}
\begin{column}{0.54\textwidth}
\begin{lstlisting}[basicstyle=\ttfamily\tiny,
  emph={Actor, pingy, pongy, runner, Logging, Props, ourSystem,terminate,sleep,actorOf,receive,become,stop}]
class Runner extends Actor {
  val log = Logging(context.system, this)
  val pingy = context.actorOf(Props[Pingy], "pingy")
  def receive = {
    case "start" =>
      val pongySys = "akka.tcp://PongyDimension@127.0.0.1:24321"
      val pongyPath = "/user/pongy"
      val url = pongySys + pongyPath
      val selection = context.actorSelection(url)
      selection ! Identify(0)
    case ActorIdentity(0, Some(ref)) =>
      pingy ! ref
    case ActorIdentity(0, None) =>
      log.info("Something's wrong - ain't no pongy anywhere!")
      context.stop(self)
    case "pong" =>
      log.info("got a pong from another dimension.")
      context.stop(self)
  }
}
\end{lstlisting}
\end{column}
\end{columns}
\end{frame}

\begin{frame}[t]\frametitle{Running the multi-dimensional Pingy-Pongy}

\begin{itemize}
  \item Start the \cod{RemotingPongySystem}
  \item Start the \cod{RemotingPingySystem} within 15 sec.
  \item Use different SBT instances
  \item Runner in \cod{PingyDimension} should get a ``pong'' soon
\end{itemize}
    
\pause
\begin{block}{Deployment logic vs. Application logic}
  \begin{itemize}
    \item Deployment log.: setting up network communication
    \item Application log.: interactions between agents
    \item These should be kept in separate
    \item In our example, \cod{Runner} handles deployment logic
  \end{itemize}
\end{block}
\end{frame}

\begin{frame}[c]\frametitle{Wrapping up remote actors}
  \begin{exampleblock}{Steps for handling remote actors}    
    \begin{itemize}
      \item \alert{Declaring} each actor system with appropriate remoting configuration
      \item \alert{Starting} each actor system in separate processes or on separate machines
      \item \alert{Obtain actor references} by using actor path selection
      \item \alert{Transparently send messages} by using these actor references
    \end{itemize}
  \end{exampleblock}
\end{frame}


\begin{frame}[t]\frametitle{Wrapping up Actors}
    % We used the Akka framework for actors for:
    \begin{itemize}
      \item Declare \alert{actor classes} and create \alert{actor instances}
      \item Model \alert{actor state} and complex \alert{actor behaviours}
      \item Manipulate the \alert{actor hierarchy} and the \alert{actor lifecycle}
      \item Use some message-passing patterns used in \alert{actor communication}
      \item Use error recovery with the built-in \alert{actor supervision} mechanism
      \item Use \alert{remote actors} to build concurrent and distributed programs
    \end{itemize}

  \splittwo{0.4}{0.6}{
    ~\\~\\
    Documentation: \url{https://doc.akka.io/docs/akka}
  }{
    \includegraphicsframed[scale=1]{src/img/actor-com.pdf}
  }
\end{frame}

\end{document}
